\section{A Reusable Trust Service}
\label{sec:service}

Rather than ship a monolithic chatbot, Le Rapporteur exposes its trust layer as
a service that others can call. The design goal is \emph{reuse}: the value is
not only the answers our own model produces, but the ability to fact-check the
citations in \emph{any} model's output.

\subsection{HTTP API}
Two endpoints cover the two use cases.
\begin{itemize}
  \item \texttt{POST /answer} --- ask a question and receive a sourced answer
  or a refusal. Each citation is annotated with \{\,\texttt{exists},
  \texttt{url}, \texttt{excerpt}\,\}; under a live backend the response also
  carries serving metrics (latency, tokens/s).
  \item \texttt{POST /verify} --- submit \emph{arbitrary} text (e.g., another
  LLM's output) and receive, per citation, whether it resolves to a real
  in-force article, plus aggregate flags such as \texttt{n\_invented} and a
  boolean \texttt{trustworthy}.
\end{itemize}
\texttt{/verify} is the key integration point: a citizen platform, a newsroom,
or a parliamentary tool can screen a third-party model's output for invented
legal citations without adopting our model at all.

\begin{lstlisting}[style=lr]
POST /verify
{ "text": "See article L. 4321-7 of the Code du travail." }

200 OK
{ "citations": [
    { "ref": "L.4321-7", "code": "Code du travail",
      "exists": false, "url": null } ],
  "n_invented": 1,
  "trustworthy": false }
\end{lstlisting}

\subsection{An MCP provider, not just a consumer}
The same two capabilities---\texttt{repondre\_question} and
\texttt{verifier\_article}---are also exposed as a \textbf{Model Context
Protocol} (MCP) server~\cite{mcp2024}, speaking JSON-RPC~2.0 over streamable
HTTP (\texttt{initialize}, \texttt{tools/list}, \texttt{tools/call}, with
\texttt{structuredContent} in the result). This inverts the usual framing: Le
Rapporteur is a \emph{provider} of verified legal knowledge, so that a
general-purpose assistant can call it \emph{as a tool} mid-conversation and
inherit the fail-closed guarantee.

\begin{lstlisting}[style=lr]
--> tools/call  verifier_article
    { "ref": "L.3121-27", "code": "Code du travail" }
<-- structuredContent
    { "exists": true, "etat": "VIGUEUR",
      "url": "https://www.legifrance.gouv.fr/.../LEGIARTI...",
      "excerpt": "La duree legale de travail effectif..." }
\end{lstlisting}

We verified interoperability end-to-end by driving our own MCP client against
this server, confirming that the tool contract---names, argument shapes, and
structured results---is honoured in both directions.

\subsection{Tool discovery and a full exchange}
An MCP client first discovers the available tools, then calls one; the server
returns \texttt{structuredContent} that the client can act on programmatically
rather than having to parse prose.

\begin{lstlisting}[style=lr]
--> tools/list
<-- { "tools": [
      { "name": "repondre_question",
        "inputSchema": { "question": "string" } },
      { "name": "verifier_article",
        "inputSchema": { "ref": "string",
                         "code": "string" } } ] }
\end{lstlisting}

A sourced \texttt{/answer} response carries exactly the provenance the verifier
produced, so a caller never has to trust the prose alone:

\begin{lstlisting}[style=lr]
POST /answer  { "question": "Duree legale du travail ?" }
200 OK
{ "answer": "La duree legale est de 35 heures ...",
  "citations": [
    { "ref": "L.3121-27", "code": "Code du travail",
      "exists": true, "etat": "VIGUEUR",
      "url": "https://www.legifrance.gouv.fr/.../LEGIARTI...",
      "excerpt": "La duree legale de travail effectif..." } ],
  "trustworthy": true,
  "metrics": { "latency_s": ..., "tok_s": ... } }
\end{lstlisting}

\subsection{Deployment modes}
The service runs in two modes behind the \emph{same} API. In \emph{live} mode a
FastAPI process talks to the model endpoint and the database. In \emph{static}
mode the pipeline's outputs for a fixed question set are precomputed by the
\emph{real} Python pipeline and served as flat files, so a zero-backend host
can present faithful, verified answers with clickable L\'egifrance links for a
public demonstration (\S\ref{sec:demo}). Because the static build is emitted by
the same code, what the public sees is what the pipeline actually
produced---not a hand-written mock.
